\chapter{A2: swapping, copy-on-write, page-table bits}
\label{ch:a2}

A2 tasks change what a page-table entry means and where a page's content lives. Three ideas
carry all of them: \textbf{a non-present entry can carry information}, \textbf{take the page
away before you touch its content}, and \textbf{a physical page may have more than one owner}.

\section{A page-table entry, three meanings}

\begin{center}
\small
\begin{tabular}{L{3.6cm}L{4.4cm}L{6.8cm}}
\toprule
\textbf{State} & \textbf{Bits} & \textbf{Meaning} \\
\midrule
not mapped & all zero & no page here (lazy: \code{loadPage} on the first access) \\
present & \code{present = 1}, \code{page\_ppn} = frame & normal; \code{writeable}, NX, \code{dirty}, \code{accessed} apply \\
swapped out (A2-01) & \code{present = 0}, marker in \code{ignored\_2}, \code{page\_ppn} = slot & the CPU ignores all bits of a non-present entry --- they are yours \\
present + clean copy (A2-05) & \code{present = 1}, \code{ignored\_1} = slot + 1 & content also on disk; \code{dirty} tells whether it is still equal \\
shared zero page (A2-09) & \code{present = 1}, \code{writeable = 0}, \code{page\_ppn} = zero page & the first write copies \\
\bottomrule
\end{tabular}
\end{center}

\begin{trapbox}
\code{ArchMemory::unmapPage} frees a page table as soon as no entry in it is \emph{present}. With
swapped entries that is wrong: it would throw away the slot numbers. Treat ``all entries zero'' as
empty (A2-01). Likewise \code{\textasciitilde ArchMemory} frees every \emph{present} page: shared
pages (A1-10) and the zero page (A2-09) must be excluded, swap slots (A2-01) freed separately.
\end{trapbox}

\section{Swap out and in --- the order is the whole task}

\begin{center}
\begin{tikzpicture}[
  box/.style={draw, rounded corners=2pt, fill=codebg, align=left, font=\footnotesize,
              inner sep=4pt, text width=6.2cm},
  arr/.style={-{Stealth[length=2mm]}, thick}]
  \node[box] (a) {\textbf{swapOut(vpn)} --- page-table lock held throughout};
  \node[box, below=2mm of a] (b) {1. \code{present = 0}, flush the TLB: nobody can write any more};
  \node[box, below=2mm of b] (c) {2. write the frame to a slot (via the ident address) --- the lock stays held:
     a thread that touches the page faults and \emph{waits}};
  \node[box, below=2mm of c] (d) {3. entry = marker + slot, free the frame};
  \draw[arr] (a) -- (b); \draw[arr] (b) -- (c); \draw[arr] (c) -- (d);
  \node[box, right=6mm of a] (e) {\textbf{page fault on a swapped page}};
  \node[box, below=2mm of e] (f) {page-table lock: still swapped? (another thread may have swapped it in already)};
  \node[box, below=2mm of f] (g) {allocate a frame, read the slot into it, free (or keep, A2-05) the slot};
  \node[box, below=2mm of g] (h) {entry = frame, marker off, \code{present = 1} \textbf{last}};
  \draw[arr] (e) -- (f); \draw[arr] (f) -- (g); \draw[arr] (g) -- (h);
\end{tikzpicture}
\end{center}

\begin{lockbox}
\begin{itemize}
\item \textbf{Disk I/O while holding the page-table Mutex is allowed} --- and here it is what keeps
  other threads off the half-written page. A SpinLock would be wrong: the request sleeps.
\item One lock order for all of A2: page-table lock $\to$ PageManager $\to$ swap bookkeeping
  (\code{SwapManager}). Your team's swap thread must use the same order.
\item Read the \code{dirty} bit only after \code{present = 0} + flush --- before that a write can still set it.
\item Removing rights (present, writable, NX) $\Rightarrow$ TLB flush in the critical section.
  Changing the frame of a present mapping (COW, relocation) $\Rightarrow$ flush.
\end{itemize}
\end{lockbox}

\section{The swap device}

\code{idea2} is a 50\,MiB partition of type 0x82 that base SWEB never uses:
\begin{lstlisting}[style=cpp]
BDVirtualDevice* swap = BDManager::getInstance()->getDeviceByName("idea2");
\end{lstlisting} Slot $s$ = bytes
$[s \cdot 4096, (s+1) \cdot 4096)$, written and read through the frame's ident address
(\code{readData}/\code{writeData}, ex08). Create the bookkeeping object in \code{startup()} after
\code{BDManager::doDeviceDetection()}; do disk I/O only from threads (the request yields until the
IDE interrupt) --- \code{startup()} itself is too early (A2-11 reads its record in
\code{ProcessRegistry::Run()}). \code{run\_test.sh} boots with a throw-away disk; \code{--keep-disk}
keeps one copy over several boots.

\section{Copy-on-write}

A write to a present read-only page faults with \code{present = 1, writing = 1} --- such faults are
rejected by \code{checkPageFaultIsValid}, so the COW branch goes \emph{before} it. CR0.WP is set
(\code{boot.32.C}), so the \textbf{kernel's} writes into user memory fault as well: a COW branch that
requires \code{user} breaks every syscall that writes into a fresh buffer. The copy is
\code{allocPPN} + \code{memcpy} between ident addresses; the zero-page case (A2-09) needs no
\code{memcpy}. Fork's COW (your mandatory part) adds reference counts per frame and two processes'
page tables --- lock both in a fixed order (e.g.\ by address) or copy the entries under the parent's
lock only.

\section{Moving pages behind a process's back}

Swap-out, relocation (A2-10), defragmentation, ``swap out all pages'' (F3): the process may run at any
moment (preemption!) and touch the page. The pattern is always the same: \emph{present = 0 first,
work, present = 1 last}, under the page-table lock; and the page-fault handler, finding a non-present
page, must first take that lock and check whether the page is back --- \emph{before} treating it as a
new page (\code{loadPage} would kill the process for a stack page). A thread that works on
\emph{other} processes' memory needs a list of address spaces with its own lock, taken
\emph{before} any page-table lock, and must keep each process alive while it works on it.

\section{Sketches for the tasks not in the course}

\begin{description}
\item[Defragmentation (F7)] Walk all swapped entries of all processes (list of address spaces as in
  A2-10), move slot contents so the used slots are contiguous, update each entry --- under each
  process' page-table lock, one slot at a time; the slot bitmap under its own lock.
\item[Deduplication of arbitrary pages] Hash present pages (a kernel thread), compare equal hashes
  byte by byte, map one frame read-only into both with a reference count, COW on write (A2-09).
\item[mmap a file] A region per process; the page fault reads the file part through
  \code{VfsSyscall::read} (like \code{Loader::readFromBinary}); \code{msync}/\code{munmap} write dirty pages back.
\item[ASLR] Randomise the stack top in \code{UserProcess} (and the heap start) with \code{rdtsc};
  the code segment needs position-independent binaries.
\item[Checksum check] Keep a checksum per swapped page (in a table or in the entry's free bits),
  verify after swap-in, kill the process (or panic) on mismatch.
\item[Unswappable pages / pages of the own process may not be swapped] A flag in the entry's free
  bits that the swap policy checks; \code{mlock}-like syscall to set it.
\end{description}
